% ============================================================================
% Fichier  : partie5/P05_C17_reconstitution.tex : FORN (LPI342)
% Sources : ISO/IEC 27042 ; NIST SP 800-86. Opération TSANAGA.
% ============================================================================

\chapter{Corrélation et Reconstitution Chronologique}
\label{chap:reconstitution}
\chaptermark{Reconstitution chronologique}

\epigraph{%
    « Une preuve isolée se discute~;\\
    un faisceau convergent s'impose. »%
}{Principe du cours}

\niveauChapitre{Expert}{Toute la Partie~IV~; Chapitre~\ref{chap:anti}.}

\begin{boxobjectifs}
\begin{itemize}
    \item Construire une chronologie à partir de plusieurs domaines.
    \item Recaler les horloges et traiter les incohérences.
    \item Faire converger les indices en un faisceau.
    \item Distinguer le certain, le probable, l'indéterminé.
    \item Préparer la démonstration pour le rapport.
\end{itemize}
\end{boxobjectifs}

\bigskip

% ============================================================================
\section*{Question directrice : que s'est-il exactement passé, et avec quelle certitude ?}
% ============================================================================

Disque, mémoire, réseau, mobile, nuage ont chacun livré des éléments. Les
assembler en une chronologie cohérente, en pesant la certitude de chacun, est
le cœur de la reconstitution. Ce chapitre l'enseigne.

\medskip
\textbf{Plan du chapitre.}
\begin{enumerate}
    \item Construire la chronologie (section~\ref{sec:c17-chrono}).
    \item Recaler et traiter les incohérences (section~\ref{sec:c17-recaler}).
    \item Le faisceau convergent (section~\ref{sec:c17-faisceau}).
    \item Certain, probable, indéterminé (section~\ref{sec:c17-certitude}).
\end{enumerate}

\separateur

% ============================================================================
\section{Construire la Chronologie}
\label{sec:c17-chrono}
% ============================================================================

\begin{boxprocedure}
\textbf{Construire une chronologie multi-domaines}
\begin{enumerate}
    \item Rassembler les événements datés de tous les domaines.
    \item Recaler les horloges (écarts connus).
    \item Ordonner et repérer les enchaînements.
    \item Noter la source et la certitude de chaque événement.
    \item Dater chaque étape de l'affaire.
\end{enumerate}
\end{boxprocedure}

% ============================================================================
\section{Recaler et Traiter les Incohérences}
\label{sec:c17-recaler}
% ============================================================================

Les horloges diffèrent (fuseau, dérive, modification volontaire). On les
recale sur une référence, on documente les écarts, et l'on traite les
incohérences~: une date impossible (un fichier modifié après son effacement)
signale souvent une manipulation (Chapitre~\ref{chap:anti}).

% ============================================================================
\section{Le Faisceau Convergent}
\label{sec:c17-faisceau}
% ============================================================================

\begin{boxscenario}{Opération TSANAGA. Le faisceau}
\textbf{Constatation.} La chronologie assemble~: 3~janvier, ouverture et copie
de \texttt{clients.xlsx} vers un amovible (disque)~; 5~janvier, téléversement
vers un service d'hébergement (réseau, trace de synchronisation)~; janvier,
messages de négociation (mobile)~; à la saisie, le fichier dans un conteneur
chiffré ouvert (mémoire).

\medskip
\textbf{Analyse} (confiance élevée). Quatre domaines convergent sur la même
histoire~: copie, envoi, négociation, détention. Aucun indice isolé ne
suffirait~; ensemble, ils forment un faisceau solide. Un second conteneur
reste NEK.
\end{boxscenario}

% ============================================================================
\section{Certain, Probable, Indéterminé}
\label{sec:c17-certitude}
% ============================================================================

\begin{boxCRO}
\textbf{Peser chaque élément.} La reconstitution distingue le \emph{certain}
(établi par plusieurs sources), le \emph{probable} (une source, à corroborer)
et l'\emph{indéterminé} (NEK). Attribuer la bonne certitude à chaque
affirmation sert la \cro{0}{+}{+}~: un rapport qui surévalue sa certitude
perd en opposabilité.
\end{boxCRO}

\begin{boxmetier}{Analyste forensique et expert judiciaire}
\textbf{Assembler sans forcer.} L'expert fait converger les domaines, mais ne
force pas le faisceau~: il laisse les NEK comme tels. Un faisceau honnête, qui
assume ses zones d'ombre, convainc davantage qu'une démonstration trop lisse.
\end{boxmetier}

\separateur

\begin{boxresume}
\textbf{Ce qu'il faut retenir de ce chapitre~:}
\begin{itemize}
    \item \textbf{La chronologie} assemble les événements datés des domaines.
    \item \textbf{On recale les horloges} et on traite les incohérences.
    \item \textbf{Le faisceau convergent} fait la force de la preuve.
    \item \textbf{On pèse} certain, probable, indéterminé.
    \item \textbf{On n'efface pas les NEK} pour lisser la démonstration.
\end{itemize}
\end{boxresume}

\begin{boxinfo}
\textbf{Points de contrôle de l'élève.} Je sais~: construire une chronologie
multi-domaines~; recaler les horloges~; bâtir un faisceau~; peser la
certitude~; préserver les NEK.
\end{boxinfo}

\bigskip

\begin{boxexo}[title={\ding{45}\quad Exercice~17.1. Chronologie \niveaubase}]
Quelles étapes pour construire une chronologie multi-domaines~?
\end{boxexo}

\begin{boxexo}[title={\ding{45}\quad Exercice~17.2. Recaler \niveaubase}]
Pourquoi recaler les horloges~?
\end{boxexo}

\begin{boxexo}[title={\ding{45}\quad Exercice~17.3. Faisceau \niveaubase}]
Pourquoi un indice isolé est-il plus fragile qu'un faisceau~?
\end{boxexo}

\begin{boxexo}[title={\ding{45}\quad Exercice~17.4. Incohérence \niveaumoyen}]
Que signale un fichier « modifié après son effacement »~?
\end{boxexo}

\begin{boxexo}[title={\ding{45}\quad Exercice~17.5. TSANAGA \niveaumoyen}]
Dressez le faisceau convergent de l'affaire, domaine par domaine.
\end{boxexo}

\begin{boxexo}[title={\ding{45}\quad Exercice~17.6. Certitude \niveaumoyen}]
Classez en certain, probable, indéterminé trois éléments de l'affaire.
\end{boxexo}

\begin{boxexo}[title={\ding{45}\quad Exercice~17.7. Surévaluer \niveaumoyen}]
Pourquoi surévaluer sa certitude nuit-il au rapport~?
\end{boxexo}

\begin{boxexo}[title={\ding{45}\quad Exercice~17.8. NEK \niveaumoyen}]
Pourquoi ne pas effacer les NEK de la reconstitution~?
\end{boxexo}

\begin{boxexo}[title={\ding{45}\quad Exercice~17.9. Chronologie complète \niveauexpert}]
Rédigez la chronologie complète et datée de l'opération TSANAGA.
\end{boxexo}

\begin{boxexo}[title={\ding{45}\quad Exercice~17.10. Carte de certitude \niveauexpert}]
Construisez la carte « élément / source / certitude » de l'affaire.
\end{boxexo}

\bigskip

\begin{boxinfo}
\textbf{Pour aller plus loin.}
\begin{itemize}
    \item ISO/IEC 27042~\cite{iso27042}~; NIST SP~800-86~\cite{nist80086}.
\end{itemize}
\end{boxinfo}

\begin{boxinfo}
\textbf{Transition.} Le faisceau établi, il faut le restituer~: la Partie~VI
s'ouvre sur le rapport forensique (Chapitre~\ref{chap:rapport}).
\end{boxinfo}
